% Storyline 3 (final), paragraph 3: information remedies, the remaining problem of judgement, the gap.
Most proposed remedies act on what observers know. Public institutions, gossip and the synchronization of opinions bring observers to a common view\cite{radzvilavicius2021,kawakatsu2024,murase2024synchronization,tham2026}, and the aggregation of several observations gives each observer more evidence to judge from\cite{michelmata2024}. Better information reduces disputed judgements, but it cannot remove the errors that make every judgement fallible\cite{hilbe2018,murase2023}. What remains is a problem of judgement, namely how a standing verdict about another person should be revised as new evidence arrives. Any revision rule must balance two errors, condemning cooperators on disputed evidence and excusing defectors whose misconduct is not yet conclusive. Human institutions often strike this balance through procedure. They escalate sanctions only with repeated violations\cite{ostrom1990,cox2010} and restore an offender's standing only after a period of good conduct\cite{blumstein2009}, and people likewise revise their impressions of others asymmetrically\cite{skowronski1989,slovic1993,mendesiedlecki2013diagnostic,siegel2018}. Models of indirect reciprocity, by contrast, mostly revise reputations through a single symmetric rule, whether the latest assessment, a score or a fixed window of observations\cite{ohtsuki2004,hilbe2018,schmid2023,michelmata2024}. How such revision rules shape the punishment dilemma, and how they should be designed, remains open.
